\begin{TheoremBox}{Verschieben, Strecken und Spiegeln}{real_functions_transformations}
  Für reelle Parameter entsteht aus einem Punkt $(u,f(u))$ des
  ursprünglichen Graphen bei der jeweiligen
  \emph{Transformation} \EnglishTerm{transformation}:
  \begin{center}
    \renewcommand{\arraystretch}{1.5}
    \begin{tabular}{@{}lll@{}}
      \toprule
      Neuer Term & Neuer Punkt & Wirkung\\ \midrule
      $f(x)+d$ & $(u,f(u)+d)$ & um $d$ nach oben\\
      $f(x-c)$ & $(u+c,f(u))$ & um $c$ nach rechts\\
      $a f(x)$ & $(u,a f(u))$ & vertikaler Faktor $a$\\
      $f(bx)$, $b\ne0$ & $(u/b,f(u))$ & horizontaler Faktor $1/b$\\
      \bottomrule
    \end{tabular}
  \end{center}
  Negative Verschiebungen wirken in die Gegenrichtung.
  Für $|a|>1$ wird vertikal gestreckt, für $0<|a|<1$ gestaucht.
  Für $|b|>1$ wird horizontal gestaucht, für $0<|b|<1$ gestreckt.
  Ein negativer Faktor bewirkt zusätzlich eine Spiegelung:
  $a<0$ an der $x$-Achse, $b<0$ an der $y$-Achse.
\end{TheoremBox}
\begin{center}
  \CachedTikzFigure{content/resources/chapters/real_functions/tikz/functions_transformations.tex}
\end{center}
\begin{KeyPointBox}{Außen direkt -- innen über die Eingabe}{real_functions_inside_outside}
  Bei $f(x)+d$ wird zum Funktionswert $d$ addiert.
  Bei $f(x-c)$ muss dagegen $x=u+c$ eingesetzt werden, damit im Inneren
  wieder die alte Eingabe $u$ steht.

  Daher verschiebt $f(x-2)$ den Graphen nach \emph{rechts}, nicht nach links.
\end{KeyPointBox}

\subsection{Mehrere Transformationen kombinieren}
\label{subsec:real_functions_combined_transformations}
Für $g(x)=a\,f(b(x-c))+d$ mit $b\ne0$ bleibt die Zuordnung der Punkte
auch bei mehreren Änderungen eindeutig:
\[
  (u,f(u))\quad\longmapsto\quad
  \left(c+\frac{u}{b},\,d+a f(u)\right),
  \qquad
  D_g=\left\{c+\frac{u}{b}\,\middle|\,u\in D_f\right\}.
\]
\begin{ExampleBox}{Eine Parabel aus bekannten Punkten konstruieren}{real_functions_transformed_parabola}
  Aus $f(x)=x^2$ soll $g(x)=-2(x-1)^2+3$ entstehen.
  Bestimmen Sie die Bilder der drei Punkte $(-1,1)$, $(0,0)$ und $(1,1)$.
  \begin{ClozeArea}[34]{7}
    Wir verschieben um $1$ nach rechts, multiplizieren die Höhen mit $-2$
    und verschieben um $3$ nach oben:
    \[
      \begin{array}{c@{\quad\longmapsto\quad}c}
        (-1,1)&(0,1)\\[0.4em]
        (0,0)&(1,3)\\[0.4em]
        (1,1)&(2,1).
      \end{array}
    \]
    Die neue Parabel ist nach unten geöffnet und zur Geraden $x=1$
    symmetrisch. Ihr höchster Punkt ist $(1,3)$.
  \end{ClozeArea}
\end{ExampleBox}
\begin{RemarkBox}
  Vor dem Ablesen einer horizontalen Verschiebung muss der innere Term
  passend ausgeklammert werden:
  $f(2x-4)=f(2(x-2))$ hat die Verschiebung $c=2$, nicht $c=4$.
\end{RemarkBox}
